\chapter*{新版前言}
\addcontentsline{toc}{chapter}{新版前言} % 手动加入目录
\markboth{新版前言}{新版前言} % 设置页眉

眼前这部《黑格尔法哲学研究》于1969年面世，旋即出了第2版以及各种译本，至70年代中期即告售罄。现在，对人们经常提到的对于新版的期盼，我再也不能置之不理了，盖因我相信，它今天又有了新的读者。因为所有这些主题：在实定法与自然法、经济与政治、市民社会与国家、现状与历史之间的近代二分以及黑格尔对其所做的概念上的消解，至今分量丝毫未减。同时，《法哲学原理》还透露了黑格尔逻辑学的一些秘密。它揭露了辩证法的历史核心。

笔者在60年代初构思这个主题时，最初只是想对黑格尔的《法哲学原理》与康德的《道德形而上学》（1797年）和费希特的《自然法基础》（1798年）进行一个结构上的比较。[在黑格尔这里发生的]基本政治概念（如“市民”、“市民社会”、“国家”）的意义断裂不容忽视。这种断裂使我从康德出发，经克里斯蒂安·沃尔夫(Christian Wolff)、莱布尼茨和霍布斯回到经院哲学的亚里士多德传统，最终回到亚里士多德的《政治学》。正是概念与其历史之间解释上的差距，向我展示了黑格尔政治哲学的另一基础。唯有如此，方有可能证明，他并不是那种导致（从R. 海谋到K. R. 波普尔）对《法哲学原理》错误阅读的威权主义的非理性的无底深渊。《法哲学原理》的另一基础就是近代革命的自然法原则，即在公民生活的风俗、伦常和伦理的传统原则之后出现的、包含在基督教之中的所有人的自由及其在法中的制度化。

为了表明这两个原则在黑格尔的以历史为中心的辩证法中的共同作用，现将书名改为《在传统与革命之间》。无须对这些研究进行改动，对此我甚感满意。第一部分第二篇（“制度中的辩证法”）是新加的。它是我于1976年在帕多瓦大学所做的一次演讲，迄今为止，只有意大利文文本（Quaderni di verifiche 2）。第二部分第三篇（“自由法则与自然的统治”）和第四部分第八篇（“黑格尔历史哲学中的进步与辩证法”）取自（同样售罄的）论文集《体系与历史》（1972年）。它们也是60年代晚期为了演讲目的而写的。

\begin{flushright}
M. 里德尔\\
1982年1月\\
于埃朗根-纽伦堡
\end{flushright}